
\subsection{滤过 $A_\infty$ 结构的 bulk 形变}
\label{subsec:3-8-5}

可以利用定理 3.8.32 定义
$C_{1+}(L;\Lambda_{0,\mathrm{nov}})$ 上 $A_\infty$ 结构
$\mathfrak m$ 的一种形变。群
$H^2(M;\Lambda^+_{0,\mathrm{nov}})$ 的元素诱导辛结构的一种扩张形变，
本小节定义的正是用于这种辛结构形变的形变；在这个意义下，可称之为
\emph{bulk 形变}\liftIndex{I010}。

令
\[
 \mathfrak b\in C^2(M;\Lambda^+_{0,\mathrm{nov}})
 =C(M;\Lambda^+_{0,\mathrm{nov}})[2]^0,
 \qquad
 \mathfrak b\equiv0\pmod{\Lambda^+_{0,\mathrm{nov}}},
\]
并置
\[
 e^{\mathfrak b}
 =
 \sum_{k=0}^{\infty}
 \underbrace{\mathfrak b\otimes\cdots\otimes\mathfrak b}_{k\text{ 个}}
 \in\widehat EC(M;\Lambda_{0,\mathrm{nov}})[2].
\]
注意（见引理 3.6.2）
\begin{equation}
 \Delta(e^{\mathfrak b})
 =e^{\mathfrak b}\otimes e^{\mathfrak b}.
\tag{3.8.37}\label{eq:3-8-37}
\end{equation}
还要注意，尽管 $E_kC[2]$ 是由 $B_kC[2]$ 对称化所得，却有
\[
 e^{\mathfrak b}
 \ne
 \sum_{k=1}^{\infty}\frac{1}{k!}\mathfrak b^{\otimes k}.
\]
这由 \eqref{eq:3-8-37} 说明。事实上，若采用另一种约定定义
$EC[2]$ 上的余积 $\Delta$，则在 $e^{\mathfrak b}$ 的定义中需要系数
$1/k!$。这里采取的办法是在定义 $\mathfrak q$ 的
(3.8.68) 中加入 $1/\ell!$。

\begin{definition}{3.8.38}\label{def:3-8-38}
设
\[
 \mathfrak b\in C^2(M;\Lambda^+_{0,\mathrm{nov}}),
 \quad
 \mathfrak b\equiv0\pmod{\Lambda^+_{0,\mathrm{nov}}},
 \qquad
 b\in C^1_{1+}(L;\Lambda^+_{0,\mathrm{nov}}).
\]
对 $k=0,1,2,\ldots$，定义
\[
\begin{aligned}
 \mathfrak m_k^{\mathfrak b,b}(x_1,\ldots,x_k)
 &:=
 \mathfrak q\bigl(
 e^{\mathfrak b}\otimes
 e^b x_1e^b\cdots e^b x_ke^b
 \bigr),\\
 \mathfrak m_k^{+\mathfrak b,b}(x_1,\ldots,x_k)
 &:=
 \mathfrak q^+\bigl(
 e^{\mathfrak b}\otimes
 e^b x_1e^b\cdots e^b x_ke^b
 \bigr).
\end{aligned}
\]
\end{definition}

\begin{lemma}{3.8.39}\label{lem:3-8-39}
设 $\delta_M\mathfrak b=0$。则
$\mathfrak m_k^{\mathfrak b,b}$ 在
$C_{1+}(L;\Lambda_{0,\mathrm{nov}})$ 上定义一个滤过
$A_\infty$ 代数结构；而且
$\mathfrak m_0^{\mathfrak b,b}=0$，当且仅当
\[
 \mathfrak q(e^{\mathfrak b}\otimes e^b)=0.
\]
类似地，$\mathfrak m_k^{+\mathfrak b,b}$ 在
$C_{1+}(L;\Lambda_{0,\mathrm{nov}})^+$ 上定义一个含幺滤过
$A_\infty$ 代数结构；而
$\mathfrak m_0^{+\mathfrak b,b}=0$，当且仅当
\[
 \mathfrak q^+(e^{\mathfrak b}\otimes e^b)=0.
\]
\end{lemma}

\begin{proof}
由假设，$\widehat\delta_Me^{\mathfrak b}=0$。把
\eqref{eq:3-8-33} 与 (3.8.34) 合并，得到
\[
\begin{aligned}
0={}&
\sum_{0\leq i\leq j\leq k}
(-1)^{\deg' x_1+\cdots+\deg' x_i}\\
&\mathfrak q\Bigl(
 e^{\mathfrak b}\otimes
 e^b x_1e^b\cdots e^b x_ie^b\,
 \mathfrak q\bigl(
  e^{\mathfrak b}\otimes
  e^b x_{i+1}e^b\cdots e^b x_je^b
 \bigr)\\
&\hspace{34mm}{}
 e^b x_{j+1}e^b\cdots e^b x_ke^b
\Bigr).
\end{aligned}
\]
这恰好是 $\mathfrak m_k^{\mathfrak b,b}$ 的 $A_\infty$ 关系。
$\mathfrak m_k^{+\mathfrak b,b}$ 的版本同理。
\end{proof}

现在研究方程 $\mathfrak m_0^{\mathfrak b,b}=0$ 有解的障碍。

\begin{definition}{3.8.40}\label{def:3-8-40}
\begin{enumerate}[label=\textup{(\arabic*)},leftmargin=2.8em]
\item 若 $\mathfrak m_0^{\mathfrak b,b}=0$ 有解，就称 $L$
\emph{经 bulk 形变后无障碍}。把所有满足
$\mathfrak m_0^{\mathfrak b,b}=0$ 的 $(\mathfrak b,b)$ 组成的集合记作
$\widehat{\mathcal M}_{\mathrm{def}}(L)$。对固定的
$\delta_M\mathfrak b=0$，置
\[
 \widehat{\mathcal M}^{\mathfrak b}(L)
 =
 \left\{
 b\in C^1(L;\Lambda^+_{0,\mathrm{nov}})
 \ \middle|\
 (\mathfrak b,b)\in
 \widehat{\mathcal M}_{\mathrm{def}}(L)
 \right\}.
\]

\item 若对某个次数为零的
$c\in\Lambda^{(0)}_{0,\mathrm{nov}}$，方程
\[
 \mathfrak m_0^{+\mathfrak b,b}(1)
 =ce\mathbf e^+
\]
有解，就称 $L$\emph{经 bulk 形变后弱无障碍}。这里
$\mathfrak m^+$ 在 $C_{1+}(L;\Lambda_{0,\mathrm{nov}})^+$ 上定义滤过
$A_\infty$ 结构，$\mathbf e^+$ 是其幺元。置
\[
 \widehat{\mathcal M}_{\mathrm{weak,def}}(L)
 =
 \left\{
 (\mathfrak b,b)
 \ \middle|\
 \mathfrak m_0^{+\mathfrak b,b}=ce\mathbf e^+
 \right\}.
\]
集合 $\widehat{\mathcal M}^{\mathfrak b}_{\mathrm{weak}}(L)$
也以类似方式定义。

\item 再定义势函数
\[
 \mathfrak{PO}\colon
 \widehat{\mathcal M}_{\mathrm{weak,def}}(L)
 \longrightarrow\Lambda^{(0)}_{0,\mathrm{nov}}
\]
使
\[
 \mathfrak m_0^{+\mathfrak b,b}(1)
 =
 \mathfrak{PO}(\mathfrak b,b)e\mathbf e^+.
\]
\end{enumerate}
\end{definition}
\liftIndex{I071}\liftIndex{I072}\liftIndex{I077}\liftIndex{I079}\liftIndex{I354}

引理 3.6.36 的类似结论将在后面给出；精确陈述见定理 4.6.52 和推论
4.6.53。还要注意，
$\widehat{\mathcal M}^{\mathfrak b}(L)$ 正是滤过 $A_\infty$ 代数
\[
 \bigl(C(L;\Lambda_{0,\mathrm{nov}}),
 \mathfrak m^{(\mathfrak b,0)}\bigr)
\]
的定界余链集合，即
\[
 \widehat{\mathcal M}^{\mathfrak b}(L)
 =
 \widehat{\mathcal M}\bigl(
 C(L;\Lambda_{0,\mathrm{nov}}),
 \mathfrak m^{(\mathfrak b,0)}
 \bigr).
\]

下面给出定理 3.6.18 的类似结论。回顾 $G(L)$ 是由
\[
 G_0(L)
 =
 \left\{
 \bigl(\omega[\beta],\mu_L(\beta)\bigr)
 \ \middle|\
 \beta\in\Pi(L),\ 
 \mathcal M(L;\beta)\ne\varnothing
 \right\}
\]
生成的幺半群。令 $G\supset G(L)$ 是满足条件 3.1.6 的幺半群，并枚举为
\[
 G=\{(\lambda_i,\mu_i)\mid i=0,1,2,\ldots\},
\]
其中
\[
 \lambda_i\leq\lambda_{i+1},
 \qquad
 \lambda_i=\lambda_{i+1}\Longrightarrow\mu_i<\mu_{i+1}.
\]

\begin{theorem}{3.8.41}\label{thm:3-8-41}
存在一列障碍类
\[
 [o_k(L;\mathrm{def})]
 \in
 \frac{H^{2-\mu_k}(L;\mathbb Q)}
 {\operatorname{Im}\bigl(
 i^*\colon H^{2-\mu_k}(M;\mathbb Q)
 \longrightarrow H^{2-\mu_k}(L;\mathbb Q)
 \bigr)},
\]
其中 $k=1,2,\ldots$；以及
\[
 b_k\in C^{1-\mu_k}(L;\mathbb Q),
 \qquad
 \mathfrak b_k\in C^{2-\mu_k}(M;\mathbb Q),
 \qquad
 \delta_M\mathfrak b_k=0,
\]
其中 $\mu_k$ 是偶数马斯洛夫指标。它们具有下列性质。

\smallskip
\noindent\textup{(3.8.42.1)}
若所有满足 $\lambda_j<\lambda_k$ 的 $j$ 所对应的
$b_j,\mathfrak b_j,o_j$ 都已定义，则闭上链
$o_k(L;\mathrm{def})$ 有定义；其上同调类
$[o_k(L;\mathrm{def})]$ 依赖于这些
$b_j,\mathfrak b_j,o_j$。

\noindent\textup{(3.8.42.2)}
若 $[o_k(L;\mathrm{def})]=0$，则
$b_k,\mathfrak b_k$ 有定义。

\noindent\textup{(3.8.42.3)}
若所有障碍类 $[o_k(L;\mathrm{def})]$ 都有定义且均为零，则
$L$ 经 bulk 形变后无障碍。
\end{theorem}

由定理 3.8.41 立即得到下列推论。

\begin{corollary}{3.8.43}\label{cor:3-8-43}
若对每个 $k$，
\[
 i^*\colon H^{2k}(M;\mathbb Q)
 \longrightarrow H^{2k}(L;\mathbb Q)
\]
都是满射，则 $L$ 经 bulk 形变后无障碍。
\end{corollary}

\begin{proof}[定理 3.8.41 的证明]
对 $k$ 作归纳。与定理 3.8.11 的证明一样，为简化记号，先假设对每个
$i$ 都有 $\lambda_i<\lambda_{i+1}$；其他情形作显然修改即可处理，见
第 3.6.2 小节。现在精确陈述归纳假设。置
\[
 b_{(i)}
 =\sum_{j=1}^{i}T^{\lambda_j}e^{\mu_j/2}b_j,
 \qquad
 \mathfrak b_{(i)}
 =\sum_{j=1}^{i}T^{\lambda_j}e^{\mu_j/2}\mathfrak b_j,
\]
并假设 $i<k$ 时的 $b_i,\mathfrak b_i$ 满足
$\delta_M\mathfrak b_i=0$，而且
\begin{equation}
 \mathfrak q\bigl(
 e^{\mathfrak b_{(i)}}\otimes e^{b_{(i)}}
 \bigr)
 \equiv0
 \pmod{T^{\lambda_i}\Lambda^+_{0,\mathrm{nov}}}
 \qquad(i<k).
\tag{3.8.44}\label{eq:3-8-44}
\end{equation}
在此假设下，将构造一个确为闭上链的
$o_k(L;\mathrm{def})$。此外，若
$[o_k(L;\mathrm{def})]$ 消失，则将证明
$b_k,\mathfrak b_k$ 有定义，并满足
\[
 \mathfrak q\bigl(
 e^{\mathfrak b_{(k)}}\otimes e^{b_{(k)}}
 \bigr)
 \equiv0
 \pmod{T^{\lambda_k}\Lambda^+_{0,\mathrm{nov}}}.
\]

现在以等式
\begin{equation}
 T^{\lambda_k}e^{\mu_k/2}o_k(L;\mathrm{def})
 \equiv
 \mathfrak q\bigl(
 e^{\mathfrak b_{(k-1)}}\otimes e^{b_{(k-1)}}
 \bigr)
 \pmod{T^{\lambda_k}\Lambda^+_{0,\mathrm{nov}}}
\tag{3.8.45}\label{eq:3-8-45}
\end{equation}
定义 $o_k(L;\mathrm{def})$。

\begin{lemma}{3.8.46}\label{lem:3-8-46}
$o_k(L;\mathrm{def})$ 是 $C(L;\mathbb Q)$ 中的闭上链。
\end{lemma}

\begin{proof}
把 \eqref{eq:3-8-33} 用于
$e^{\mathfrak b_{(k-1)}}\otimes e^{b_{(k-1)}}$，得到
\[
\begin{aligned}
 \mathfrak q\Bigl(
 e^{\mathfrak b_{(k-1)}}\otimes
 e^{b_{(k-1)}}\,
 \mathfrak q\bigl(
  e^{\mathfrak b_{(k-1)}}\otimes e^{b_{(k-1)}}
 \bigr)
 e^{b_{(k-1)}}
 \Bigr)=0.
\end{aligned}
\]
由 \eqref{eq:3-8-45}，
\[
\begin{aligned}
T^{\lambda_k}e^{\mu_k/2}
\mathfrak q\Bigl(
 e^{\mathfrak b_{(k-1)}}\otimes
 e^{b_{(k-1)}}o_k(L;\mathrm{def})e^{b_{(k-1)}}
\Bigr)
\equiv0
\pmod{T^{\lambda_k}\Lambda^+_{0,\mathrm{nov}}}.
\end{aligned}
\]
利用 \eqref{eq:3-8-34-2} 以及
\[
 b_{(k-1)}\equiv\mathfrak b_{(k-1)}
 \equiv0\pmod{\Lambda^+_{0,\mathrm{nov}}},
\]
便有
$\overline{\mathfrak m}_1(o_k(L;\mathrm{def}))=0$。
\translatornote{原书此处写
$b_{(k)}\equiv\mathfrak b_{(k)}\equiv0$；然而在本归纳步骤中
$b_k,\mathfrak b_k$ 尚未构造，紧前两式实际代入的也都是第
$k-1$ 阶部分和。译文依归纳型别恢复为
$b_{(k-1)},\mathfrak b_{(k-1)}$。}
\end{proof}

把 $o_k(L;\mathrm{def})$ 诱导的类定义为
\[
 [o_k(L;\mathrm{def})]
 \in
 \frac{H^{2-\mu_k}(L;\mathbb Q)}
 {\operatorname{Im}\bigl(
 i^*\colon H^{2-\mu_k}(M;\mathbb Q)
 \longrightarrow H^{2-\mu_k}(L;\mathbb Q)
 \bigr)}.
\]
现在设 $[o_k(L;\mathrm{def})]=0$。这意味着存在
\[
 b_k\in C^{1-\mu_k}(L;\mathbb Q),
 \qquad
 \mathfrak b_k\in C^{2-\mu_k}(M;\mathbb Q)
\]
使
\translatornote{原书这里再次把第一项写成
$C^{1-\mu_k}_2(L;\mathbb Q)$；但定理 3.8.41 的正式陈述不带下标
$2$，全章也未定义与此处相容的 $C_2(L)$。译文沿用定理 3.1.15
处已经说明的裁决，删去这一排字残片。}
\begin{equation}
 i^\pitchfork\mathfrak b_k
 +\overline{\mathfrak m}_1(b_k)
 +o_k(L;\mathrm{def})=0,
 \qquad
 \delta_M(\mathfrak b_k)=0.
\tag{3.8.47}\label{eq:3-8-47}
\end{equation}
计算
$\mathfrak q(e^{\mathfrak b_{(k)}}\otimes e^{b_{(k)}})$，得到
\begin{equation}
\begin{aligned}
&\mathfrak q\bigl(
 e^{\mathfrak b_{(k)}}\otimes e^{b_{(k)}}
 \bigr)
 -
 \mathfrak q\bigl(
 e^{\mathfrak b_{(k-1)}}\otimes e^{b_{(k-1)}}
 \bigr)\\
&\quad\equiv
T^{\lambda_k}e^{\mu_k/2}
\left(
 \mathfrak q_{1,0}(\mathfrak b_k\otimes1)
 +\mathfrak q_{0,1}(1\otimes b_k)
\right)
\pmod{T^{\lambda_k}\Lambda^+_{0,\mathrm{nov}}}.
\end{aligned}
\tag{3.8.48}\label{eq:3-8-48}
\end{equation}
由 \eqref{eq:3-8-34-1} 与 \eqref{eq:3-8-34-2}，右端模
$T^{\lambda_k}\Lambda^+_{0,\mathrm{nov}}$ 等于
\[
 T^{\lambda_k}e^{\mu_k/2}
 \bigl(
 i^\pitchfork\mathfrak b_k
 +\overline{\mathfrak m}_1b_k
 \bigr).
\]
因此
\begin{equation}
 \mathfrak q\bigl(
 e^{\mathfrak b_{(k)}}\otimes e^{b_{(k)}}
 \bigr)
 \equiv0
 \pmod{T^{\lambda_k}\Lambda^+_{0,\mathrm{nov}}},
\tag{3.8.49}\label{eq:3-8-49}
\end{equation}
如所要求。这就证明了 (3.8.42.1) 与 (3.8.42.2)；而
(3.8.42.3) 的证明与定理 3.6.18 中 (3.6.19.4) 的证明相同。
\end{proof}
下面把定理 3.8.41 推广到经 bulk 形变后弱无障碍的情形。这可视为定理
3.6.43 的类似版本。仍把幺半群 $G\supset G(L)$ 枚举为
\[
 G=\{(\lambda_i,\mu_i)\mid i=0,1,2,\ldots\},
\]
其中
\[
 \lambda_i\leq\lambda_{i+1},
 \qquad
 \lambda_i=\lambda_{i+1}\Longrightarrow\mu_i<\mu_{i+1}.
\]
如第 3.6.3 小节，置
\[
 G_{\mu\leq0}
 =
 \{(\lambda_i,\mu_i)\in G\mid\mu_i\leq0\}.
\]
从给定的 $G$ 枚举中删去所有满足 $\mu_i>0$ 的 $i$，就得到一个
非递减函数
\[
 j\colon\mathbb Z_{\geq0}\longrightarrow\mathbb Z_{\geq0},
 \qquad j(0)=0,
\]
使
\[
 G_{\mu\leq0}
 =
 \{(\lambda_{j(i)},\mu_{j(i)})\mid i=0,1,2,\ldots\}.
\]
见 \eqref{eq:3-6-42}。在我们的几何情形中，假设 3.6.39 显然成立。
还要指出：忽略 bulk 形变后，定理 3.6.43 的证明立即化为下面给出的
证明。

\begin{theorem}{3.8.50}\label{thm:3-8-50}
存在一列障碍类
\[
 [o_k(L;\mathrm{weak,def})]
 \in
 \frac{H^{2-\mu_k}(L;\mathbb Q)}
 {\operatorname{Im}\bigl(
 i^*\colon H^{2-\mu_k}(M;\mathbb Q)
 \longrightarrow H^{2-\mu_k}(L;\mathbb Q)
 \bigr)},
\]
其中 $k=1,2,\ldots$；以及
\[
 b_k\in C^{1-\mu_k}(L;\mathbb Q),
 \qquad
 \mathfrak b_k\in C^{2-\mu_k}(M;\mathbb Q),
 \qquad
 \delta_M\mathfrak b_k=0,
\]
其中 $\mu_k$ 是偶数马斯洛夫指标。它们具有下列性质。

\smallskip
\noindent\textup{(3.8.51.1)}
若 $i<k$ 时的 $o_i(L;\mathrm{weak,def})$、
$b_i,\mathfrak b_i$ 均已定义，并且 $\mu_k\leq0$，则
$[o_k(L;\mathrm{weak,def})]$ 有定义；它依赖于
$i<k$ 时的 $o_i(L;\mathrm{weak,def})$、
$b_i,\mathfrak b_i$。

\noindent\textup{(3.8.51.2)}
当 $\mu_k\leq0$ 时，若
$[o_k(L;\mathrm{weak,def})]$ 已定义且为零，则
$b_k,\mathfrak b_k$ 有定义。当 $\mu_k>0$ 时，若
$i<k$ 时的 $b_i,\mathfrak b_i$ 均存在，则
$b_k,\mathfrak b_k$ 有定义。

\noindent\textup{(3.8.51.3)}
若所有满足 $\mu_k\leq0$ 的障碍类
$[o_k(L;\mathrm{weak,def})]$ 都有定义且均为零，则
$L$ 经 bulk 形变后弱无障碍。
\end{theorem}
\liftIndex{I108}\liftIndex{I349}

\begin{proof}
证明方法与定理 3.8.41 相同。仍为简化记号，假设对每个 $i$ 都有
$\lambda_i<\lambda_{i+1}$，并置
\[
 \mathfrak b_{(i)}
 =\sum_{j=1}^{i}T^{\lambda_j}e^{\mu_j/2}\mathfrak b_j,
 \qquad
 b_{(i)}
 =\sum_{j=1}^{i}T^{\lambda_j}e^{\mu_j/2}b_j.
\]
若每个 $\ell\leq k$ 都属于 $\operatorname{Im}(j)$，则对
$[o_k(L;\mathrm{weak,def})]$ 与 $b_k,\mathfrak b_k$ 的论证和定理
3.8.41 完全相同，归纳便进入下一步。主要问题是下面的情形；其他情形
也可类似处理。假设每个 $\ell<k$ 都属于 $\operatorname{Im}(j)$，而
$k\notin\operatorname{Im}(j)$。由于 $\mu_k$ 是偶数，故
$\mu_k\geq2$。

回顾 $\mathbf e=\operatorname{PD}[L]$ 是滤过 $A_\infty$ 代数的同伦
幺元，并如第 3.3 节考虑
\[
\begin{aligned}
 C_{1+}(L;\Lambda_{0,\mathrm{nov}})^+
 ={}&C_{1+}(L;\Lambda_{0,\mathrm{nov}})\\
 &{}\oplus\Lambda_{0,\mathrm{nov}}\mathbf e^+
 \oplus\Lambda_{0,\mathrm{nov}}\mathbf f,
\end{aligned}
\]
其中
\[
 \deg\mathbf e^+=0,\qquad\deg\mathbf f=-1.
\]
由 \eqref{eq:3-3-5-2}，
\[
 \overline{\mathfrak m}_1^+(\mathbf f)
 \equiv\mathbf e^+-\mathbf e
 \pmod{\Lambda^+_{0,\mathrm{nov}}},
\]
所以
\[
 H(L;\mathbb Q)
 \cong
 H\bigl(C(L;\mathbb Q)^+,\overline{\mathfrak m}_1^+\bigr).
\]
定理 3.8.32 还把算子 $\mathfrak q$ 延拓为 $\mathfrak q^+$。

在上述情形下，假设 $i<k$ 时已有
\[
 b_i\in C_{1+}(L;\mathbb Q)^+,
 \qquad
 \mathfrak b_i\in C(M;\mathbb Q),
 \qquad
 \delta_M\mathfrak b_i=0,
\]
并且
\[
 \mathfrak q^+\bigl(
 e^{\mathfrak b_{(i)}}\otimes e^{b_{(i)}}
 \bigr)
 \equiv c_i e\mathbf e^+
 \pmod{T^{\lambda_i}\Lambda^+_{0,\mathrm{nov}}}
 \qquad(i<k),
\]
其中 $c_i\in\Lambda^{(0)}_{0,\mathrm{nov}}$。
\translatornote{原书仅在上式中把“$(\mathrm{weak})$”分别置于两个
指数之后，既未定义相应的新部分和，紧接着的第 $k$ 步也恢复使用
$\mathfrak b_{(k)},b_{(k)}$。译文依本证明开头的定义统一删去这两个
游离标记。}
在这个假设下，将构造障碍上链
$o_k(L;\mathrm{weak,def})$ 以及 $b_k,\mathfrak b_k$，使
$o_k(L;\mathrm{weak,def})$ 是闭上链，并且
\[
 \mathfrak q^+\bigl(
 e^{\mathfrak b_{(k)}}\otimes e^{b_{(k)}}
 \bigr)
 \equiv c_k e\mathbf e^+
 \pmod{T^{\lambda_k}\Lambda^+_{0,\mathrm{nov}}}
\]
对某个 $c_k\in\Lambda^{(0)}_{0,\mathrm{nov}}$ 成立。

像前面一样，以等式
\[
 T^{\lambda_k}e^{\mu_k/2}o_k(L;\mathrm{weak,def})
 \equiv
 \mathfrak q^+\bigl(
 e^{\mathfrak b_{(k-1)}}\otimes e^{b_{(k-1)}}
 \bigr)
 \pmod{T^{\lambda_k}\Lambda^+_{0,\mathrm{nov}}}
\]
定义 $o_k(L;\mathrm{weak,def})$。同理可证：

\begin{lemma}{3.8.52}\label{lem:3-8-52}
$o_k(L;\mathrm{weak,def})$ 是
$C^{2-\mu_k}(L;\mathbb Q)^+$ 中的
$\overline{\mathfrak m}_1^+$-闭上链。
\end{lemma}

\begin{proof}
若 $\mathbf e=\operatorname{PD}[L]$ 是严格幺元，就不必考虑
$C(L;\Lambda_{0,\mathrm{nov}})^+$、$\mathfrak q^+$ 等扩张对象，此时
可用与引理 3.8.46 相同的论证证明结论。

若 $\operatorname{PD}[L]$ 是同伦幺元，则定义 3.8.38 与引理
3.8.39 表明：对每个满足
\[
 \mathfrak b\in C^2(M;\Lambda^+_{0,\mathrm{nov}}),
 \qquad \delta_M\mathfrak b=0
\]
的 $\mathfrak b$，
\[
 \bigl(C(L;\Lambda_{0,\mathrm{nov}})^+,
 \mathfrak m^{+\mathfrak b}\bigr)
\]
是含幺滤过 $A_\infty$ 代数。于是仍可像严格幺元情形那样证明结论。
\translatornote{原书此处引用“定义 3.8.37 与引理 3.8.38”，但
(3.8.37) 是余积公式，定义与引理实际依次编号为 3.8.38、3.8.39；
译文据紧前对象改正。}
\end{proof}

先考虑 $\mu_k>2$。由于假设 $L$ 可定向，故 $\mu_k\geq4$。
由次数理由，$o_k(L;\mathrm{weak,def})$ 自动上同调于零；也就是说，
它是一个同调维数大于 $\dim L$ 的闭链。这里还需注意，
$\deg\mathbf f=-1$。此外，即使 $\mu_k=3$，次数为 $-1$ 的
$C(L;\mathbb Q)^+$ 元素在 $\mathbb Q$ 上也不可能是闭上链，因为由
\eqref{eq:3-3-5-2}，
\[
 \overline{\mathfrak m}_1^+(\mathbf f)
 \equiv\mathbf e^+-\mathbf e\ne0
 \pmod{\Lambda^+_{0,\mathrm{nov}}}.
\]
置 $\mathfrak b_k=0$，并可找到 $b_k$ 使
\[
 \delta b_k=o_k(L;\mathrm{weak,def}).
\]
于是显然有
\[
 \mathfrak q^+\bigl(
 e^{\mathfrak b_{(k)}}\otimes e^{b_{(k)}}
 \bigr)
 \equiv c_ke\mathbf e^+
 \pmod{T^{\lambda_k}\Lambda^+_{0,\mathrm{nov}}},
\]
其中 $c_k=c_{k-1}$，如所要求。

若 $\mu_k=2$，则可写成
\[
\begin{aligned}
o_k(L;\mathrm{weak,def})
&=\overline c_{k,1}\operatorname{PD}[L]
  +\overline c_{k,2}\mathbf e^+\\
&\in C^0(L;\mathbb Q)^+
 =C^0(L;\mathbb Q)\oplus\mathbb Q\mathbf e^+,
\end{aligned}
\]
其中 $\overline c_{k,i}\in\mathbb Q$。在 $\mu_k=2$ 的情形下，并不
假设
\[
 [o_k(L;\mathrm{weak,def})]
 =0
 \in
 \frac{H^0(L;\mathbb Q)}
 {\operatorname{Im}\bigl(
 i^*\colon H^0(M;\mathbb Q)\longrightarrow H^0(L;\mathbb Q)
 \bigr)}.
\]
定义
\begin{align}
 b_k&=\overline c_{k,1}\mathbf f,
\tag{3.8.53.1}\label{eq:3-8-53-1}\\
 \mathfrak b_k&=0.
\tag{3.8.53.2}\label{eq:3-8-53-2}
\end{align}
注意 \eqref{eq:3-3-5-2} 与 \eqref{eq:3-8-35-4}。与定理 3.8.41
证明相同的论证给出
\translatornote{原书在下式第三项把分量写成
$\mathfrak q_{0,1}$，但输入 $b_k=\overline c_{k,1}\mathbf f$ 属于含
同伦幺元的扩张复形；而本句所引 (3.8.35.4) 正是
$\mathfrak q^+_{0,1}(1\otimes\mathbf f)$ 的公式。译文据此补回上标
$+$。}
\[
\begin{aligned}
&\mathfrak q^+\bigl(
 e^{\mathfrak b_{(k)}}\otimes e^{b_{(k)}}
 \bigr)
 -
 \mathfrak q^+\bigl(
 e^{\mathfrak b_{(k-1)}}\otimes e^{b_{(k-1)}}
 \bigr)\\
&\quad\equiv
 \mathfrak q^+_{0,1}
 \bigl(1\otimes T^{\lambda_k}e^{\mu_k/2}b_k\bigr)\\
&\quad=
 \overline c_{k,1}(\mathbf e^+-\mathbf e)
 T^{\lambda_k}e^{\mu_k/2}
 \pmod{T^{\lambda_k}\Lambda^+_{0,\mathrm{nov}}},
\end{aligned}
\]
这是因为
\[
 \mathfrak h_2^{+M}
 \bigl(1,1\mathbin{\overline\otimes}1\bigr)
 \equiv0\pmod{\Lambda^+_{0,\mathrm{nov}}}.
\]
因此
\[
\begin{aligned}
 \mathfrak q^+\bigl(
 e^{\mathfrak b_{(k)}}\otimes e^{b_{(k)}}
 \bigr)
&\equiv
o_k(L;\mathrm{weak,def})T^{\lambda_k}e^{\mu_k/2}\\
&\quad+
\overline c_{k,1}(\mathbf e^+-\mathbf e)
T^{\lambda_k}e^{\mu_k/2}\\
&\equiv
(\overline c_{k,1}+\overline c_{k,2})
T^{\lambda_k}e\mathbf e^+\\
&=c_ke\mathbf e^+
\pmod{T^{\lambda_k}\Lambda^+_{0,\mathrm{nov}}},
\end{aligned}
\]
其中
\[
 c_k=(\overline c_{k,1}+\overline c_{k,2})T^{\lambda_k}.
\]
故归纳成立。另一部分同理。事实上，置
\[
 \mathfrak b=\lim_{k\to\infty}\mathfrak b_{(k)},
 \qquad
 b=\lim_{k\to\infty}b_{(k)},
 \qquad
 c=\lim_{k\to\infty}c_k,
\]
则这些极限关于能量滤过在完备化中收敛，并满足
\[
 \mathfrak m_0^{+\mathfrak b,b}=ce\mathbf e^+.
\]
定理 3.8.50 证毕。
\end{proof}

下面利用 (3.8.36) 证明命题 3.8.54。为简化记号，本段以
$\mathfrak q$ 和 $C(L;\Lambda_{0,\mathrm{nov}})^+$ 分别代替
$\mathfrak q^{++}$ 和
$C_{1++}(L;\Lambda_{0,\mathrm{nov}})^+$。

\begin{proposition}{3.8.54}\label{prop:3-8-54}
设\translatornote{原书把 $\mathfrak b$ 所在的环境链复形写成
$C^2(L;\Lambda^+_{0,\mathrm{nov}})^+$；但
$\mathfrak b'=\mathfrak b-ce\widetilde{\mathbf e}^+$ 使用的是
$\operatorname{PD}[M]$ 所给的环境幺元，且
$\mathfrak q^{++}$ 的第一输入属于 $C(M)^+$。译文依型别恢复为
$M$。同一段引导句还把定理 3.8.32 定义的
$C_{1++}(L;\Lambda_{0,\mathrm{nov}})^+$ 漏写了下标 $1$；上文一并
恢复。}
\[
 \mathfrak b\in C^2(M;\Lambda^+_{0,\mathrm{nov}})^+,
 \qquad
 c\in\Lambda^+_{0,\mathrm{nov}},
 \qquad
 \deg c=0.
\]
置
\[
 \mathfrak b'
 =\mathfrak b-ce\widetilde{\mathbf e}^+.
\]
则有：
\begin{align}
 \widehat{\mathcal M}^{\mathfrak b}_{\mathrm{weak}}(L)
 &=
 \widehat{\mathcal M}^{\mathfrak b'}_{\mathrm{weak}}(L),
\tag{3.8.55.1}\label{eq:3-8-55-1}\\
 \mathfrak m_k^{\mathfrak b',b}
 &=
 \mathfrak m_k^{\mathfrak b,b}
 \qquad(k\ne0),
\tag{3.8.55.2}\label{eq:3-8-55-2}\\
 \mathfrak m_0^{\mathfrak b',b}(1)
 &=
 \mathfrak m_0^{\mathfrak b,b}(1)-ce\mathbf e^+.
\tag{3.8.55.3}\label{eq:3-8-55-3}
\end{align}
这里 \eqref{eq:3-8-55-2} 对上述集合中的任意 $b$ 成立。特别地，
\[
 \mathfrak{PO}(\mathfrak b',b)
 =
 \mathfrak{PO}(\mathfrak b,b)-ce.
\]
\end{proposition}

\begin{proof}
令 $\mathbf x\in B_kC(L;\Lambda_0)$，其中 $k\ne0$。利用
\eqref{eq:3-8-36-1} 与 \eqref{eq:3-8-36-2}，计算得
\begin{equation}
\begin{aligned}
 \mathfrak q(e^{\mathfrak b'},\mathbf x)
 &=
 \mathfrak q(e^{\mathfrak b},\mathbf x)
 -ce\,\mathfrak q
   (e^{\mathfrak b}\widetilde{\mathbf e}^+,\mathbf x)
 +\cdots\\
 &=
 \mathfrak q(e^{\mathfrak b},\mathbf x).
\end{aligned}
\tag{3.8.56.1}\label{eq:3-8-56-1}
\end{equation}
另一方面，
\begin{equation}
\begin{aligned}
\mathfrak q(e^{\mathfrak b'},1)
={}&
\mathfrak q(e^{\mathfrak b},1)
-ce\,\mathfrak q(\widetilde{\mathbf e}^+,1)\\
&{}-ce\,\mathfrak q
\bigl((e^{\mathfrak b}-1)\widetilde{\mathbf e}^+,1\bigr)\\
&{}+c^2e^2\mathfrak q
\bigl(e^{\mathfrak b}(\widetilde{\mathbf e}^+)^2,1\bigr)
+\cdots\\
={}&
\mathfrak q(e^{\mathfrak b},1)-ce\mathbf e^+.
\end{aligned}
\tag{3.8.56.2}\label{eq:3-8-56-2}
\end{equation}
(3.8.55) 很容易由 (3.8.56) 推出。
\end{proof}

\begin{corollary}{3.8.57}\label{cor:3-8-57}
若
$b\in\widehat{\mathcal M}^{\mathfrak b}_{\mathrm{weak}}(L)$，
则存在某个 $\mathfrak b'$，使
$b\in\widehat{\mathcal M}^{\mathfrak b'}(L)$。
\end{corollary}

\begin{proof}
只须取
\[
 \mathfrak b'
 =
 \mathfrak b-\mathfrak{PO}(\mathfrak b,b)e
 \widetilde{\mathbf e}^+.
\]
\end{proof}

现在假设 $L$ 经 bulk 形变后无障碍，并令
\[
 (\mathfrak b,b),(\mathfrak b,b')
 \in\widehat{\mathcal M}_{\mathrm{def}}(L).
\]
注意两对的第一分量 $\mathfrak b$ 相同，第二分量则可以不同；也就是说，
\[
 b,b'\in\widehat{\mathcal M}^{\mathfrak b}(L).
\]
下面定义上链复形
\[
 CF\bigl((L;(\mathfrak b,b)),(L;(\mathfrak b,b'))\bigr).
\]

\begin{definition}{3.8.58}\label{def:3-8-58}
\begin{enumerate}[label=\textup{(\arabic*)},leftmargin=2.8em]
\item 当 $L$ 经 bulk 形变后无障碍时，作为分次模置
\[
 CF\bigl((L;(\mathfrak b,b)),(L;(\mathfrak b,b'))\bigr)
 =
 C(L;\Lambda_{0,\mathrm{nov}}).
\]
其微分\translatornote{原书在本定义两段引导句的第二个微分下标中均
漏写 $b'$ 的撇号，而紧随其后的 (3.8.59.1)、(3.8.59.2) 都正确使用
$b'$；译文据定义的两个输入统一补回。}定义为
\begin{equation}
\begin{aligned}
 \delta_{(\mathfrak b,b),(\mathfrak b,b')}(x)
 &=
 \mathfrak m_1^{\mathfrak b,b,b'}(x)\\
 &:=
 \mathfrak q\bigl(
 e^{\mathfrak b},e^bxe^{b'}
 \bigr).
\end{aligned}
\tag{3.8.59.1}\label{eq:3-8-59-1}
\end{equation}

\item 当 $L$ 经 bulk 形变后弱无障碍时，如第 3.3 节置
\[
 CF\bigl((L;(\mathfrak b,b)),(L;(\mathfrak b,b'))\bigr)
 =
 C(L;\Lambda_{0,\mathrm{nov}}),
\]
并定义其微分
$\delta^+_{(\mathfrak b,b),(\mathfrak b,b')}$ 为
\begin{equation}
\begin{aligned}
 \delta^+_{(\mathfrak b,b),(\mathfrak b,b')}(x)
 &=
 \mathfrak m_1^{+\mathfrak b,b,b'}(x)\\
 &:=
 \mathfrak q^+\bigl(
 e^{\mathfrak b},e^bxe^{b'}
 \bigr).
\end{aligned}
\tag{3.8.59.2}\label{eq:3-8-59-2}
\end{equation}
在不致混淆时，有时也把它简写为
$\delta_{(\mathfrak b,b),(\mathfrak b,b')}$。按构造，右端属于
\[
 C(L;\Lambda_{0,\mathrm{nov}})
 \subset C(L;\Lambda_{0,\mathrm{nov}})^+,
\]
故可应用注 3.6.34 (2)。
\end{enumerate}
\end{definition}

当两种情形中都有 $b=b'$ 时，分别简写为
\[
\begin{aligned}
 \delta_{(\mathfrak b,b)}(x)
 &=
 \mathfrak m_1^{\mathfrak b,b}(x)
 =
 \mathfrak q(e^{\mathfrak b},e^bxe^b),\\
 \delta^+_{(\mathfrak b,b)}(x)
 &=
 \mathfrak m_1^{+\mathfrak b,b}(x)
 =
 \mathfrak q^+(e^{\mathfrak b},e^bxe^b).
\end{aligned}
\]

\begin{lemma}{3.8.60}\label{lem:3-8-60}
\begin{enumerate}[label=\textup{(\arabic*)},leftmargin=2.8em]
\item 若 $L$ 经 bulk 形变后无障碍，则对
\[
 (\mathfrak b,b),(\mathfrak b,b')
 \in\widehat{\mathcal M}_{\mathrm{def}}(L)
\]
有
\[
 \delta_{(\mathfrak b,b),(\mathfrak b,b')}
 \circ
 \delta_{(\mathfrak b,b),(\mathfrak b,b')}
 =0.
\]

\item 若 $L$ 经 bulk 形变后弱无障碍，则对
\[
 (\mathfrak b,b),(\mathfrak b,b')
 \in\widehat{\mathcal M}_{\mathrm{weak,def}}(L)
\]
只要
\[
 \mathfrak{PO}(\mathfrak b,b)
 =
 \mathfrak{PO}(\mathfrak b,b'),
\]
就有
\[
 \delta^+_{(\mathfrak b,b),(\mathfrak b,b')}
 \circ
 \delta^+_{(\mathfrak b,b),(\mathfrak b,b')}
 =0.
\]
\end{enumerate}
\end{lemma}

\begin{proof}
结论由引理 3.8.39、引理 3.6.8 和注 3.6.34 (2) 得到。
\end{proof}

\begin{definition}{3.8.61}\label{def:3-8-61}
\begin{enumerate}[label=\textup{(\arabic*)},leftmargin=2.8em]
\item 对
\[
 (\mathfrak b,b),(\mathfrak b,b')
 \in\widehat{\mathcal M}_{\mathrm{def}}(L),
\]
定义由这两对数据形变的弗洛尔上同调
\[
 HF\bigl(
 (L;(\mathfrak b,b)),(L;(\mathfrak b,b'));
 \Lambda_{0,\mathrm{nov}}
 \bigr)
\]
为下列上链复形的上同调群\translatornote{原书在此等号右端把
定义 3.8.58 刚命名的复形 $CF$ 写成 $C$，并在第二个对象外多置
一层左括号；译文依该定义恢复 $CF$，并统一两边的对象分组。}：
\[
\begin{aligned}
&HF\bigl(
 (L;(\mathfrak b,b)),(L;(\mathfrak b,b'));
 \Lambda_{0,\mathrm{nov}}
 \bigr)\\
&\quad=
H\Bigl(
 CF\bigl((L;(\mathfrak b,b)),(L;(\mathfrak b,b'))\bigr),
 \delta_{(\mathfrak b,b),(\mathfrak b,b')}
\Bigr).
\end{aligned}
\]
\item 当 $L$ 经 bulk 形变后弱无障碍时，对满足
\[
 \begin{gathered}
 (\mathfrak b,b),(\mathfrak b,b')
 \in\widehat{\mathcal M}_{\mathrm{weak,def}}(L),\\
 \mathfrak{PO}(\mathfrak b,b)
 =
 \mathfrak{PO}(\mathfrak b,b')
 \end{gathered}
\]
的两对数据，以类似方式定义
\[
 HF\bigl(
 (L;(\mathfrak b,b)),(L;(\mathfrak b,b'));
 \Lambda_{0,\mathrm{nov}}
 \bigr).
\]
\end{enumerate}
在两种情形中，若 $b=b'$，都简写为
\[
 HF(L;(\mathfrak b,b);\Lambda_{0,\mathrm{nov}}).
\]
\end{definition}
\liftIndex{I106}\liftIndex{I329}

下一章将研究弗洛尔上同调如何依赖于 $(\mathfrak b,b)$；它与镜像对称
的关系已在第 1.4 节提及。

下面转向定理 3.8.32 的另一个应用。

\begin{theorem}{3.8.62}\label{thm:3-8-62}
对每个
$(\mathfrak b,b)\in\widehat{\mathcal M}_{\mathrm{def}}(L)$，存在同态
\[
\begin{aligned}
 i^*_{\mathfrak{qm},(\mathfrak b,b)}\colon
 H^*(M;\Lambda_{0,\mathrm{nov}})
 \longrightarrow
 HF\bigl(
 (L;(\mathfrak b,b)),(L;(\mathfrak b,b));
 \Lambda_{0,\mathrm{nov}}
 \bigr).
\end{aligned}
\]
这里 $H^*(M;\Lambda_{0,\mathrm{nov}})$ 是通常的
$\Lambda_{0,\mathrm{nov}}$-系数上同调群。上述同态由上链同态
\[
 i^\pitchfork_{\mathfrak{qm},(\mathfrak b,b)}
 \colon
 C(M;\Lambda_{0,\mathrm{nov}})
 \longrightarrow
 C_{1+}(L;\Lambda_{0,\mathrm{nov}})
\]
诱导，且
\begin{equation}
 i^\pitchfork_{\mathfrak{qm},(\mathfrak b,b)}
 \equiv
 \pm i^\pitchfork
 \pmod{\Lambda^+_{0,\mathrm{nov}}}.
\tag{3.8.63}\label{eq:3-8-63}
\end{equation}
当 $L$ 经 bulk 形变后弱无障碍时，同样的结论也成立。
\end{theorem}

\begin{proof}
令 $x\in C(M;\Lambda_{0,\mathrm{nov}})$，定义
\[
 i^\pitchfork_{\mathfrak{qm},(\mathfrak b,b)}(x)
 =
 (-1)^{\deg x}
 \mathfrak q\bigl(
 e^{\mathfrak b}xe^{\mathfrak b}\otimes e^b
 \bigr).
\]
把
$e^{\mathfrak b}xe^{\mathfrak b}\otimes e^b$
代入 \eqref{eq:3-8-33}，并取其
$B_1C_{1+}(L)$ 分量，得到
\[
\begin{aligned}
&(-1)^{\deg' x}
\mathfrak q\Bigl(
 e^{\mathfrak b}xe^{\mathfrak b}\otimes
 \mathfrak q(e^{\mathfrak b}\otimes e^b)
\Bigr)\\
&\quad+
\mathfrak q\bigl(
 e^{\mathfrak b}\delta_M(x)e^{\mathfrak b}\otimes e^b
\bigr)\\
&\quad+
\mathfrak q\Bigl(
 e^{\mathfrak b}\otimes
 e^b\,
 \mathfrak q(
  e^{\mathfrak b}xe^{\mathfrak b}\otimes e^b
 )e^b
\Bigr)
=0.
\end{aligned}
\]
由于
$(\mathfrak b,b)\in\widehat{\mathcal M}_{\mathrm{def}}(L)$，第一项为零。
故
\begin{equation}
\begin{aligned}
 i^\pitchfork_{\mathfrak{qm},(\mathfrak b,b)}
 \bigl(\delta_M(x)\bigr)
 =
 \pm\delta_{(\mathfrak b,b),(\mathfrak b,b)}
 \left(
 i^\pitchfork_{\mathfrak{qm},(\mathfrak b,b)}(x)
 \right).
\end{aligned}
\tag{3.8.64}\label{eq:3-8-64}
\end{equation}
这说明
$i^\pitchfork_{\mathfrak{qm},(\mathfrak b,b)}$
是上链同态。\eqref{eq:3-8-63} 由
\eqref{eq:3-8-34-1} 得到。经 bulk 形变后弱无障碍的情形，可利用
\eqref{eq:3-8-35-1} 处理。
\end{proof}

定理 3.8.62 的一部分还可以用第 6 章将要构造的谱序列重新表述。
